\section{ROS Attacks}
\label{background:ros}

Unlike prior schemes in the literature~\cite{DrijversEFKLNS19,BenhamoudaLLOR20},
both FROST in Chapter~\ref{chp:frost} and Arctic in Chapter~\ref{chp:arctic} are secure against concurrent adversaries which may attempt ROS-style attacks.
%
The ROS (Random Inhomogeneities in a Overdetermined Solvable System of Linear Equations) problem~\cite{schnorr06,BenhamoudaLLOR20,Schnorr01,FuchsbauerPS20}
over $\Zq$ in dimension $\ell$
is defined as follows.
At a high level,
the goal is to find a solution to $\ell+1$ linear equations modulo $q$,
where some elements in the equations are output from a random oracle.
In other words,
given a random oracle $\HROS$ and a set of arbitrary functions $f_1, \ldots, f_{\ell+1}$
where $f_i(\Zq^\ell) \rightarrow \Zq$,
the goal is to find coefficients $a_{k, i} \in \Zq, i \in \set{\ell}$ such that
\[
  a_{k,1} c_1 + a_{k,2} c_2 + \ldots + a_{k,\ell} c_\ell = \HROS(f_{\ell+1}((a_{k,i})_{i \in \set{\ell}}, \aux))
\]

\noindent where $c_i = \HROS(f_i((a_{k,i})_{i \in \set{\ell}}, \aux))$.
It is always possible to find the first $\ell$ values simply by setting
$a_{ki} = 1, i \in \set{\ell}$;
the challenge is to find the $\ell+1\supth$ solution.

Wagner~\cite{Wagner02} and Benhamouda et al.~\cite{BenhamoudaLLOR20} showed how to solve the ROS problem in subexponential time,
and demonstrated its applicability to the forgeability in a range of blind and multi-party
signature schemes---including blind and threshold Schnorr signatures--- assuming a \emph{concurrent adversary},
which may open an arbitrary number of signing sessions simultaneously.

